\subsection{RQ2: content signatures on public data}
\label{sec:rq2}

Table~\ref{tab:rq2-pooled} pools all labelled items per surface. The content signatures are
\emph{precise and blind}: on 37{,}058 benign prompts the prompt rules fire on 0.91\% (95\% CI 0.82 to
1.01), and on 9{,}467 benign documents and 10{,}001 real MCP tool descriptions on 0.25\% and 0.17\%,
but they recall only 11.9\% of 19{,}611 attack prompts, 0.9\% of 16{,}272 injected documents and
23.7\% of 485 poisoned tool descriptions.

\begin{table}[t]\centering\small
\resizebox{\textwidth}{!}{\begin{tabular}{llrlrl}
\toprule
Surface & Rule & Benign $n$ & FPR \% [95\% CI] & Attack $n$ & Recall \% [95\% CI] \\
\midrule
prompt & LLMD-001 & 37,058 & 0.84 [0.76, 0.94] & 19,611 & 10.6 [10.2, 11.1] \\
prompt & LLMD-011 & 37,058 & 0.04 [0.02, 0.07] & 19,611 & 0.0 [0.0, 0.1] \\
prompt & LLMD-012 & 37,058 & 0.02 [0.01, 0.05] & 19,611 & 1.6 [1.4, 1.7] \\
prompt & any rule & 37,058 & 0.91 [0.82, 1.01] & 19,611 & 11.9 [11.4, 12.3] \\
document & LLMD-002 & 9,467 & 0.00 [0.00, 0.04] & 16,272 & 0.6 [0.5, 0.8] \\
document & LLMD-021 & 9,467 & 0.25 [0.17, 0.38] & 16,272 & 0.9 [0.8, 1.1] \\
document & any rule & 9,467 & 0.25 [0.17, 0.38] & 16,272 & 0.9 [0.8, 1.1] \\
tool description & LLMD-015 & 10,001 & 0.17 [0.11, 0.27] & 485 & 23.7 [20.1, 27.7] \\
tool description & any rule & 10,001 & 0.17 [0.11, 0.27] & 485 & 23.7 [20.1, 27.7] \\
\bottomrule
\end{tabular}
}
\caption{Content signatures pooled over all labelled public datasets, per surface.}
\label{tab:rq2-pooled}
\end{table}

\paragraph{Recall tracks phrasing, not intent.} Per dataset (Table~\ref{tab:rq2-datasets}), recall is
highest on Lakera Gandalf (39.6\%), whose prompts were selected for similarity to ``Ignore all
previous instructions'', and near zero where attackers had reason to avoid that phrasing: 1.0\% on
LLMail-Inject, where participants iterated against a defended email agent, and 0.0\% on NVIDIA's agentic
indirect injections, which pose as late additions from the user (``Hi Assistant, I forgot to mention
earlier \ldots''). Role-play jailbreaks (TrustAIRLab, jackhhao: 4.5\% and 5.6\%) are largely outside
what an injection signature is meant to match.

\paragraph{Where the false positives come from.} The highest false-positive rate is on the
TrustAIRLab ``regular'' prompts (6.6\%), collected from the same prompt-sharing communities as its
jailbreaks. Inspection shows that almost all of these hits are prompt-marketplace templates that open
with ``Please ignore all previous instructions. Respond only in [TARGETLANGUAGE]'': the
instruction-override phrase is a benign idiom in that ecosystem, used to reset a chat before a
template. On unlabelled real traffic (WildChat-1M) the prompt rules alert on
WildChat-1M: 1.80\% of 20,000 prompts
a rate that a SOC would see as a steady stream of low-value alerts.

\paragraph{Implication.} Content signatures are cheap to run and their alerts are easy to explain,
which makes them useful as context and as one input to correlation; as a primary detector they
miss the attacks that matter. This motivates the question of RQ1: do the structural and behavioural
rules, which do not depend on the attacker's wording, close the gap on live attacks?
